% Transcription — fonds Alexandre Grothendieck, shelfmark 34, batch 11 (pages 201-209).
% Established from the facsimile archives/batches/34/batch-11.pdf, per the skill
% transcribe-grothendieck.
% The batch interleaves two materials. Odd pages 201, 203, 205, 207, 209 are
% leaves of an English typescript by another author (pages "2.12", "2.3", "2.5",
% "2.11", "2.10" of a text on reductive group actions, 1-PS, Iwahori's theorem);
% under RIGHTS.md they are not transcribed, only noted. Even pages 202, 204, 206,
% 208 are his own French notes on weights in l-adic cohomology: spectral
% sequences attached to a normal crossings divisor, degeneration (d_r = 0 for
% r >= 3) by a weight argument, and independence of l.
% Pass: Opus 5.5 (claude-opus-5-5), 2026-10-03 — first pass, unchecked against the pages by a human.
%
\documentclass[11pt,a4paper]{article}
\input{../preamble/grothendieck}

\folder{34}
\batch{11}
\pages{201}{209}
\foldertitle{SGA 7 : notes manuscrites (s.d.).}
\dating{[vers 1967-1973]}
\watermark{Édition de démonstration}

\begin{document}

\page{201}
\note{Dactylographie en anglais d'une autre main, non signée ni datée ; feuillet paginé « -2.12- » d'un texte sur les actions de groupes réductifs (fin d'une démonstration, Proposition 2.4 sur la propreté d'une action, théorème d'Iwahori). Non transcrite (droits de tiers).}

\page{202}
\note{la phrase commence sur un feuillet antérieur au lot : le sujet de « induit » manque ici.}
induit un hom. injectif $\operatorname{Im}\varphi^{i} \to E_{2}^{0,i}$, i.e.
$\operatorname{Im}\varphi^{i} \cap \operatorname{Filt}^{1}(H^{i}(Y,\mathbb{Q}_{\ell})) = 0$,
de sorte que \struck{$\operatorname{Coker}\varphi^{i}$} la filtration quotient sur
$\operatorname{Coker}\varphi^{i}$ a comme \struck{\ill{}}
$\operatorname{Coker}(H^{i}(X) \to E_{2}^{0,i}) =$ défaut d'exactitude de
\[
  H^{i}(X) \longrightarrow \prod H^{i}(Y_{\alpha}) \longrightarrow \prod H^{i}(Y_{\alpha\beta}) ,
\]
et les $E_{2}^{p,q}$ $(p \geq 1,\ p+q = i) =$
\struck{$\prod H^{q}(Y_{\alpha_{0}\dots\alpha_{p}})$}
$H^{p}(\sigma \mapsto H^{q}(Y_{\sigma}))$.
Cette filtration est également indépendante de la résolution choisie des
singularités, et correspond à la filtration de $H_{!}^{i+1}(X)$\note{l'exposant se lit $i+1$ ; il est peut-être $i-1$ ou $i$.}
par poids décroissants (\uncertain{fonction} des poids $i$).

\marginal{de plus, si conj. \add{\uncertain{fantaisies}} standard, les dimensions des divers
quotients de la filtration \uncertain{poids} sur $H_{!}^{i}(X,\mathbb{Q}_{\ell})$ sont
indépendantes de $\ell$ !}
\note{note écrite en oblique dans la marge gauche, à hauteur des lignes précédentes ; le mot inséré au-dessus de « conj. » est d'une lecture douteuse.}

(b) Par dualité, on peut considérer aussi
\[
  \cdots \longrightarrow H^{i}_{Y}(X') \longrightarrow H^{i}(X') \longrightarrow H^{i}(X) \longrightarrow ,
\]
\note{le dernier terme se lit $H^{i}(X)$ ; l'$X$ y porte une marque (prime ou indice) indistincte.}
on a d'autre part
\[
  H^{a}_{Y}(X) \Longleftarrow E_{2}^{pq} = H^{p}(Y, \underline{H}^{q}_{Y}(\mathbb{Q}_{\ell\,X'})) ,
\]
\struck{Or d'autre part}
$\underline{H}^{q}_{Y}(\mathbb{Q}_{\ell}) = 0$ si $q < 2$, et
\[
  \underline{H}^{q}_{Y}(\mathbb{Q}_{\ell}) = \bigwedge^{q-1}\Bigl(\coprod_{\alpha} \mathbb{Q}_{\ell\,Y_{\alpha}}(-1)\Bigr)
  = \coprod_{\alpha_{1}<\dots<\alpha_{q-1}} \mathbb{Q}_{\ell\,Y_{\alpha_{1}\dots\alpha_{q-1}}}(-1)
\]
\note{avant le dernier $\coprod$, un symbole biffé ; le twist final est écrit $(-1)$, alors que la ligne suivante porte $(-(q-1))$.}
donc
\[
  E_{2}^{pq} =
  \begin{cases}
    0 & \text{si } q < 2 \\
    \displaystyle\coprod_{\alpha_{1}<\dots<\alpha_{q-1}} H^{p}(Y_{\alpha_{1}\dots\alpha_{q-1}})(-(q-1)) & \text{si } q \geq 2
  \end{cases}
\]
donc $E_{2}^{pq}$ est de poids $p + 2q - 2$, donc on \struck{a pour} aura
\struck{$d_{r} = 0$ si}
\[
  d_{r} = 0 \quad \text{si } r \geq 3 .
\]

\page{203}
\note{Dactylographie de la même main et du même texte, feuillet « -2.3- » (Théorème 2.1, critère numérique de semi-stabilité par les sous-groupes à un paramètre). Quelques marques au crayon dans la marge (renvois de numérotation), d'attribution incertaine, ne sont pas transcrites.}

\page{204}
\struck{Il faut donc prouver que $E_{3}^{p,q} =$}

Il faut écrire que $E_{2}^{p*}$ est (\add{modulo twist}) le
\struck{\ill{}} complexe dual \add{de (a)} de
\struck{$E_{2}^{p*}$ \ill{}}
$\sigma \mapsto H^{p}(Y_{\sigma})$, compris dans le sens de la dualité de Poincaré,
de sorte que
\struck{\ill{}}
${}^{b}E_{3}^{**}$ est dual de ${}^{(a)}E_{2}^{**}$, à twist près.
\note{« (a) » et « (b) » renvoient aux deux présentations des pages précédentes ; la seconde est le (b) de la p.~202. Un passage biffé en boucle sous « $\sigma \mapsto H^{p}(Y_{\sigma})$ » n'est pas lisible.}

(c) Autre présentation \add{de $H^{*}(X,\mathbb{Q}_{\ell})$} : suite spectrale de Leray de
$g : X \hookrightarrow X'$ ; on a
\[
  H^{*}(X) \Longleftarrow E_{2}^{pq} = H^{p}(X', R^{q}g_{*}(\mathbb{Q}_{\ell}))
\]
Or
\begin{align*}
  R^{q}g_{*}(\mathbb{Q}_{\ell}) &\simeq \bigwedge^{q} \coprod_{\alpha} \mathbb{Q}_{\ell}(-1)_{Y_{\alpha}} \\
  &\simeq \coprod_{\alpha_{1}<\dots<\alpha_{q}} \mathbb{Q}_{\ell\,Y_{\alpha_{1}\dots\alpha_{q}}}(-q) .
\end{align*}
Donc
\[
  \begin{cases}
    E_{2}^{pq} = \displaystyle\coprod_{\alpha_{1}<\dots<\alpha_{q}} H^{p}(Y_{\alpha_{1}\dots\alpha_{q}})(-q) & \text{si } q \geq 1 \\
    E_{2}^{p0} = H^{p}(X')
  \end{cases}
\]
qui est \emph{de poids} $p + 2q$. \emph{Donc on trouve} \struck{\ill{}}
\[
  d_{r} = 0 \quad \text{si } r \geq 3 ,
\]
et il faut calculer encore $E_{3}^{pq}$, \struck{comme} i.e. le
\note{la suite de la phrase est le diagramme ; à droite, séparé par un trait vertical : « si $q \geq 2$, donc $q-1 \geq 1$ ». Les annotations « dim $n-1-q$ » et « dim $n-q$ » sont écrites au-dessus des deux termes de la ligne médiane, « dual » le long des flèches verticales.}

\begin{tikzcd}[column sep=small, row sep=small, nodes={font=\scriptsize}]
  E_{2}^{pq} \arrow[r] \arrow[d, no head, "{=}" description] & E_{2}^{p+2,q-1} \arrow[d, no head, "{=}" description] \\
  \coprod_{\alpha_{1}<\dots<\alpha_{q}} H^{p}(Y_{\alpha_{1}\dots\alpha_{q}})(-q) \arrow[r] \arrow[d, no head, "\text{dual}"'] & \coprod_{\alpha_{1}<\dots<\alpha_{q-1}} H^{p+2}(Y_{\alpha_{1}\dots\alpha_{q-1}})(-(q-1)) \arrow[d, no head, "\text{dual}"] \\
  H^{2n-2-p-2q}(Y_{\alpha_{1}\dots\alpha_{q}})(n-1) & H^{2n-2q-p-2}(Y_{\alpha_{1}\dots\alpha_{q-1}})(n-1)
\end{tikzcd}

\note{les exposants et les twists de la dernière ligne sont d'une lecture incertaine.}

\page{205}
\note{Dactylographie de la même main et du même texte, feuillet « -2.5- » (stabilité propre, fonctions invariantes homogènes). En marge, au crayon, « 1.11 » corrigeant un renvoi « 1.9 » du texte, d'attribution incertaine ; non transcrit.}

\page{206}
③

\begin{tikzcd}
  X \arrow[d, "{f\ \text{(propre)}}"'] & U \arrow[l, hook'] \arrow[ld, no head, "g" description] \arrow[d, no head, "f_{U}"] \\
  S & \eta \arrow[l, hook']
\end{tikzcd}

\note{croquis : $X \supset U$ au-dessus de $S \hookleftarrow \eta$, $f$ (propre) à gauche, $f_{U}$ à droite, la lettre $g$ entre les deux ; à côté de $\eta$ un mot biffé. La disposition des traits de $g$ et de $f_{U}$ est restituée approximativement.}

$X$ régulier, $U = X -{}$ \struck{$X$} $Y$, $Y = \bigcup Y_{\alpha} \supset X_{0}$,
\uncertain{en} $Y$,
$\left\{\begin{array}{l} Y_{\alpha} \text{ div. réguliers} \\ Y \text{ div. à croisements normaux.} \end{array}\right.$

[\emph{NB} faire pour simplifier $Y = X_{0}$, $U = X_{\eta}$]

\[
  H^{*}(U) \Longleftarrow E_{2}^{pq} = H^{p}(X_{0}, R^{q}g_{*}(\mathbb{Q}_{\ell\,X})|X_{0})
\]
\[
  R^{q}g_{*}(\mathbb{Q}_{\ell\,X}) \simeq \bigwedge^{q} \coprod_{\alpha} (\mathbb{Q}_{\ell}(-1)_{Y_{\alpha}})
  \simeq \coprod_{\alpha_{1}<\dots<\alpha_{q}} \mathbb{Q}_{\ell\,Y_{\alpha_{1}\dots\alpha_{q}}}(-q)
\]
donc
\[
  \begin{cases}
    E_{2}^{p,q} = \text{\struck{$H^{p}$}}\ \displaystyle\coprod_{\alpha_{1}<\dots<\alpha_{q}} H^{p}((Y_{\alpha_{1}\dots\alpha_{q}})_{0})(-q) & \text{si } q \geq 1 \\
    E_{2}^{p0} = H^{p}(X_{0})
  \end{cases}
\]
— la première de poids $p + 2q = i + q$, la seconde de poids \struck{\ill{}} $p = i$ \struck{\ill{}}.
\note{à gauche de l'accolade, une tache d'encre ; dans $H^{p}((Y_{\alpha_{1}\dots\alpha_{q}})_{0})$ un symbole biffé précède la parenthèse intérieure.}

$E_{r}^{pq} \to E_{r}^{p+r,\,q-r+1}$ :

a) $q - r + 1 \neq 0$ : poids $p+2q$ et $p + r + 2(q-r+1)$, \uncertain{nulle} si $r \neq 2$.

b) $q - r + 1 = 0$ \struck{\ill{} $p+2q$} poids $p+2q$ et $\leq p+r$ ; il faut pour que
$d_{r}^{pq} \neq 0$ que l'on ait $p+2q \leq p+r$, i.e. $r \geq 2q$ ; comme $r = q+1$,
cela s'écrit $q+1 \geq 2q$, i.e. $q \leq 1$, donc $r \leq 2$.

\emph{Donc} $\underline{d_{r} = 0}$ si $r \geq 3$.

\struck{En dimension}

$H^{i}$ filtration dont les termes, en croissant, sont

$E_{2}^{i,0}$ de poids \struck{\ill{}} $i$

$E_{2}^{i-1,1}$ de poids $i+1$

$E_{2}^{i-2,2}$ — $i+2$

$E_{2}^{0,i}$ — $2i$

[lui-même muni d'une filtration convenable qui correspond aux différents poids…]

\page{207}
\note{Dactylographie de la même main et du même texte, feuillet « -2.11- » (calcul en coordonnées via le théorème d'Iwahori). Non transcrite.}

\page{208}
\struck{Les $d_{2}$}
Modulo identification des $d_{2}$ par morphismes canoniques (de nature « motivique »),
on trouve que la dimension des $E_{\infty}^{pq}$ est indépendante de $\ell$. De plus, les
\struck{valeurs propres} pol. car. de Frobenius \add{inverse} y opérant sont à coeff.
\uncertain{entiers} \ill{} indépendants de $\ell$.

\marginal{compatibilité à vérifier…}

Utilisons la suite exacte
\[
  0 \longrightarrow H^{i-1}(X_{\bar\eta})_{I} \longrightarrow H^{i}(X_{\eta}) \longrightarrow H^{i}(X_{\bar\eta})^{I} \longrightarrow 0
  \tag{6}
\]
\note{le numéro de la suite, à gauche, se lit « (6) » ou « (b) » ; dans le terme médian un premier $X$ est biffé ; les indices $\bar\eta$ sont restitués d'après un $\eta$ surmonté d'un trait.}

Nous \struck{voulons} \add{obtiendrons} donc des informations sur
$H^{i}(X_{\bar\eta})^{I}$ et $H^{i}(X_{\bar\eta})_{I}$, savoir

a) leurs dimensions indépendantes de $\ell$ $(\neq p)$

b) Les pol. car. de Frobenius inverse opérant sur eux sont à coeff. entiers
indépendants de $\ell$, et les val. propres \ill{} \ill{} des val. absolues $q^{i/2}$
dans $H^{i}(X_{\bar\eta})^{I}$ et $H^{i}(X_{\bar\eta})_{I}$.
\note{la base de $q^{i/2}$ est d'une lecture incertaine.}

\struck{Il suffit de le prouver pour $i \leq n$, car si $i > n$ \ill{} mais par \ill{}}
\note{deux lignes biffées, reliées par un trait sinueux ; seul le début se lit.}

On procède par récurrence sur $i$. Pour $i = 0$, on a
$H^{0}(X_{\eta}) \simeq H^{0}(X_{\bar\eta})^{I}$, on gagne. Pour $i > 0$, on
\struck{$H^{i}(X_{\bar\eta})$} utilise \uncertain{(b)}, \struck{\ill{}} et l'hyp.
de récurrence…
\note{la page s'arrête sur ces points de suspension ; la suite n'est pas dans le lot.}

\page{209}
\note{Dactylographie de la même main et du même texte, feuillet « -2.10- » (énoncé et démonstration d'un critère de propreté pour l'orbite, théorème d'Iwahori). Non transcrite.}

\end{document}
